% Options for packages loaded elsewhere
\PassOptionsToPackage{unicode}{hyperref}
\PassOptionsToPackage{hyphens}{url}
\documentclass[
  11pt,
  a4paper,
]{article}
\usepackage{xcolor}
\usepackage[margin=18mm]{geometry}
\usepackage{amsmath,amssymb}
\setcounter{secnumdepth}{-\maxdimen} % remove section numbering
\usepackage{iftex}
\ifPDFTeX
  \usepackage[T1]{fontenc}
  \usepackage[utf8]{inputenc}
  \usepackage{textcomp} % provide euro and other symbols
\else % if luatex or xetex
  \usepackage{unicode-math} % this also loads fontspec
  \defaultfontfeatures{Scale=MatchLowercase}
  \defaultfontfeatures[\rmfamily]{Ligatures=TeX,Scale=1}
\fi
\ifPDFTeX\else
  % xetex/luatex font selection
    \setmainfont[]{Noto Serif}
    \setsansfont[]{Noto Sans}
    \setmonofont[]{Noto Sans Mono}
  \setmathfont[]{Noto Sans Math}
\fi
% Use upquote if available, for straight quotes in verbatim environments
\IfFileExists{upquote.sty}{\usepackage{upquote}}{}
\IfFileExists{microtype.sty}{% use microtype if available
  \usepackage[]{microtype}
  \UseMicrotypeSet[protrusion]{basicmath} % disable protrusion for tt fonts
}{}
\makeatletter
\@ifundefined{KOMAClassName}{% if non-KOMA class
  \IfFileExists{parskip.sty}{%
    \usepackage{parskip}
  }{% else
    \setlength{\parindent}{0pt}
    \setlength{\parskip}{6pt plus 2pt minus 1pt}}
}{% if KOMA class
  \KOMAoptions{parskip=half}}
\makeatother
\ifLuaTeX
\usepackage[bidi=basic]{babel}
\else
\usepackage[bidi=default]{babel}
\fi
\babelprovide[main,import]{russian}
\ifPDFTeX
\else
\babelfont{rm}[]{Noto Serif}
\fi
% get rid of language-specific shorthands (see #6817):
\let\LanguageShortHands\languageshorthands
\def\languageshorthands#1{}
\setlength{\emergencystretch}{3em} % prevent overfull lines
\providecommand{\tightlist}{%
  \setlength{\itemsep}{0pt}\setlength{\parskip}{0pt}}
\usepackage{bookmark}
\IfFileExists{xurl.sty}{\usepackage{xurl}}{} % add URL line breaks if available
\urlstyle{same}
\hypersetup{
  pdftitle={Билеты по дискретным случайным процессам},
  pdflang={ru-RU},
  hidelinks,
  pdfcreator={LaTeX via pandoc}}

\title{Билеты по дискретным случайным процессам}
\author{}
\date{}

\begin{document}
\maketitle

\section{Билет 14. Коэффициенты корреляции и евклидова геометрия
центрированных величин со вторым
моментом}\label{ux431ux438ux43bux435ux442-14.-ux43aux43eux44dux444ux444ux438ux446ux438ux435ux43dux442ux44b-ux43aux43eux440ux440ux435ux43bux44fux446ux438ux438-ux438-ux435ux432ux43aux43bux438ux434ux43eux432ux430-ux433ux435ux43eux43cux435ux442ux440ux438ux44f-ux446ux435ux43dux442ux440ux438ux440ux43eux432ux430ux43dux43dux44bux445-ux432ux435ux43bux438ux447ux438ux43d-ux441ux43e-ux432ux442ux43eux440ux44bux43c-ux43cux43eux43cux435ux43dux442ux43eux43c}

Источники: Сердобольская, лекция 6; Ширяев 1, §11.

\subsection{1. Короткий
ответ}\label{ux43aux43eux440ux43eux442ux43aux438ux439-ux43eux442ux432ux435ux442}

Если случайные величины имеют конечные вторые моменты, то их можно
рассматривать как векторы в гильбертовом пространстве
\(L^2(\Omega,\mathcal F,P)\) со скалярным произведением

\[
\langle X,Y\rangle=\mathbb E[XY].
\]

Для центрированных величин

\[
X^0=X-\mathbb EX,\qquad Y^0=Y-\mathbb EY
\]

это скалярное произведение равно ковариации:

\[
\langle X^0,Y^0\rangle
=
\mathbb E[(X-\mathbb EX)(Y-\mathbb EY)]
=
\operatorname{Cov}(X,Y).
\]

Коэффициент корреляции при \(\operatorname{Var}X>0\) и
\(\operatorname{Var}Y>0\):

\[
\rho(X,Y)
=
\frac{\operatorname{Cov}(X,Y)}
{\sqrt{\operatorname{Var}X\,\operatorname{Var}Y}}
=
\frac{\langle X^0,Y^0\rangle}{\|X^0\|_2\|Y^0\|_2}.
\]

Поэтому \(\rho(X,Y)\) - это косинус угла между центрированными векторами
\(X^0\) и \(Y^0\):

\[
\rho(X,Y)=\cos\angle(X^0,Y^0),
\qquad
|\rho(X,Y)|\le1.
\]

Равенство \(\rho=0\) означает некоррелированность, то есть
ортогональность центрированных величин в \(L^2\). Это не означает
независимость в общем случае.

\subsection{2. Полный
ответ}\label{ux43fux43eux43bux43dux44bux439-ux43eux442ux432ux435ux442}

\subsubsection{2.1. Введение и геометрическая
интерпретация}\label{ux432ux432ux435ux434ux435ux43dux438ux435-ux438-ux433ux435ux43eux43cux435ux442ux440ux438ux447ux435ux441ux43aux430ux44f-ux438ux43dux442ux435ux440ux43fux440ux435ux442ux430ux446ux438ux44f}

Этот билет продолжает Билет 12 о гильбертовых пространствах и связывает
его с вероятностными понятиями ковариации и корреляции. До этого
случайные величины рассматривались как измеримые функции, распределения
и интегралы. Здесь важна новая точка зрения: случайные величины с
конечным вторым моментом можно складывать, умножать на числа, измерять
их длины и углы между ними.

\subsubsection{\texorpdfstring{2.2. Пространство \(L^2\) и скалярное
произведение}{2.2. Пространство L\^{}2 и скалярное произведение}}\label{ux43fux440ux43eux441ux442ux440ux430ux43dux441ux442ux432ux43e-l2-ux438-ux441ux43aux430ux43bux44fux440ux43dux43eux435-ux43fux440ux43eux438ux437ux432ux435ux434ux435ux43dux438ux435}

Пусть задано вероятностное пространство

\[
(\Omega,\mathcal F,P).
\]

Пространство

\[
L^2(\Omega,\mathcal F,P)
=
\{X:\mathbb E X^2<\infty\}/\sim
\]

состоит из случайных величин с конечным вторым моментом, где величины,
равные почти наверное, считаются одним элементом. В вещественном случае
скалярное произведение задается формулой

\[
\langle X,Y\rangle=\mathbb E[XY].
\]

Оно корректно определено, потому что из неравенства Коши-Буняковского
следует

\[
\mathbb E|XY|
\le
\sqrt{\mathbb EX^2}\sqrt{\mathbb EY^2}
<\infty.
\]

Норма в \(L^2\):

\[
\|X\|_2=\sqrt{\mathbb EX^2}.
\]

Расстояние между двумя случайными величинами:

\[
\|X-Y\|_2
=
\sqrt{\mathbb E[(X-Y)^2]}.
\]

Это среднеквадратичное расстояние. Именно поэтому геометрия \(L^2\)
естественна для дисперсий, ковариаций, регрессии, прогнозирования и
корреляционных характеристик случайных процессов.

\subsubsection{2.3. Центрирование и
ковариация}\label{ux446ux435ux43dux442ux440ux438ux440ux43eux432ux430ux43dux438ux435-ux438-ux43aux43eux432ux430ux440ux438ux430ux446ux438ux44f}

Теперь выделим центрированные величины. Для любой \(X\in L^2\) положим

\[
X^0=X-\mathbb EX.
\]

Тогда

\[
\mathbb EX^0=0,
\qquad
\|X^0\|_2^2
=
\mathbb E[(X-\mathbb EX)^2]
=
\operatorname{Var}X.
\]

Для двух величин \(X,Y\in L^2\) ковариация:

\[
\operatorname{Cov}(X,Y)
=
\mathbb E[(X-\mathbb EX)(Y-\mathbb EY)].
\]

В геометрических обозначениях:

\[
\operatorname{Cov}(X,Y)=\langle X^0,Y^0\rangle.
\]

То есть ковариация - это скалярное произведение центрированных версий
\(X\) и \(Y\).

\subsubsection{2.4. Коэффициент корреляции как косинус
угла}\label{ux43aux43eux44dux444ux444ux438ux446ux438ux435ux43dux442-ux43aux43eux440ux440ux435ux43bux44fux446ux438ux438-ux43aux430ux43a-ux43aux43eux441ux438ux43dux443ux441-ux443ux433ux43bux430}

Коэффициент корреляции определяется, если обе дисперсии положительны:

\[
\operatorname{Var}X>0,
\qquad
\operatorname{Var}Y>0.
\]

Тогда

\[
\rho(X,Y)
=
\frac{\operatorname{Cov}(X,Y)}
{\sqrt{\operatorname{Var}X\,\operatorname{Var}Y}}.
\]

Через геометрию \(L^2\) это

\[
\rho(X,Y)
=
\frac{\langle X^0,Y^0\rangle}{\|X^0\|_2\|Y^0\|_2}.
\]

В обычной евклидовой геометрии для ненулевых векторов \(u,v\):

\[
\cos\angle(u,v)
=
\frac{\langle u,v\rangle}{\|u\|\|v\|}.
\]

Значит

\[
\rho(X,Y)=\cos\angle(X^0,Y^0).
\]

Именно это нужно проговорить как главный смысл билета: корреляция -
нормированная ковариация, то есть косинус угла между центрированными
случайными величинами.

\subsubsection{2.5. Неравенство Коши-Буняковского и диапазон
значений}\label{ux43dux435ux440ux430ux432ux435ux43dux441ux442ux432ux43e-ux43aux43eux448ux438-ux431ux443ux43dux44fux43aux43eux432ux441ux43aux43eux433ux43e-ux438-ux434ux438ux430ux43fux430ux437ux43eux43d-ux437ux43dux430ux447ux435ux43dux438ux439}

Из этой формулы сразу следует

\[
|\rho(X,Y)|\le1.
\]

Доказательство - это неравенство Коши-Буняковского:

\[
|\langle X^0,Y^0\rangle|
\le
\|X^0\|_2\|Y^0\|_2.
\]

Подставляя

\[
\langle X^0,Y^0\rangle=\operatorname{Cov}(X,Y),
\qquad
\|X^0\|_2^2=\operatorname{Var}X,
\qquad
\|Y^0\|_2^2=\operatorname{Var}Y,
\]

получаем

\[
|\operatorname{Cov}(X,Y)|
\le
\sqrt{\operatorname{Var}X\,\operatorname{Var}Y},
\]

а после деления на положительный знаменатель:

\[
|\rho(X,Y)|\le1.
\]

\subsubsection{2.6. Особые случаи: некоррелированность и
независимость}\label{ux43eux441ux43eux431ux44bux435-ux441ux43bux443ux447ux430ux438-ux43dux435ux43aux43eux440ux440ux435ux43bux438ux440ux43eux432ux430ux43dux43dux43eux441ux442ux44c-ux438-ux43dux435ux437ux430ux432ux438ux441ux438ux43cux43eux441ux442ux44c}

Разберем смысл особых случаев.

Если

\[
\rho(X,Y)=0,
\]

то

\[
\operatorname{Cov}(X,Y)=0
\]

и

\[
X^0\perp Y^0
\]

в \(L^2\). Такие величины называются некоррелированными. Важно:
некоррелированность в общем случае слабее независимости. Из
независимости при существовании вторых моментов действительно следует

\[
\operatorname{Cov}(X,Y)=0,
\]

потому что

\[
\mathbb E[XY]=\mathbb EX\,\mathbb EY.
\]

Но обратное неверно.

Классический пример: пусть \(X\) равномерна на \([-1,1]\), а

\[
Y=X^2.
\]

Тогда \(Y\) полностью определяется через \(X\), то есть независимости
нет. Но

\[
\operatorname{Cov}(X,Y)
=
\mathbb E[X^3]-\mathbb EX\,\mathbb EX^2
=
0-0\cdot\mathbb EX^2
=0,
\]

поскольку распределение \(X\) симметрично и \(\mathbb E X^3=0\). Значит
зависимые величины могут быть некоррелированными.

\subsubsection{2.7. Линейная зависимость и вырожденная
дисперсия}\label{ux43bux438ux43dux435ux439ux43dux430ux44f-ux437ux430ux432ux438ux441ux438ux43cux43eux441ux442ux44c-ux438-ux432ux44bux440ux43eux436ux434ux435ux43dux43dux430ux44f-ux434ux438ux441ux43fux435ux440ux441ux438ux44f}

Если

\[
|\rho(X,Y)|=1,
\]

то в неравенстве Коши-Буняковского достигается равенство. Это означает
линейную зависимость центрированных величин в \(L^2\):

\[
Y^0=aX^0
\quad\text{почти наверное}
\]

для некоторого \(a\ne0\). Эквивалентно,

\[
Y=aX+b
\quad\text{почти наверное}.
\]

Если \(a>0\), то \(\rho=1\); если \(a<0\), то \(\rho=-1\).

Если одна из дисперсий равна нулю, например

\[
\operatorname{Var}X=0,
\]

то \(X=\mathbb EX\) почти наверное, то есть \(X\) является константой в
\(L^2\). Тогда угол с ней как с центрированным вектором не определен,
потому что

\[
X^0=0.
\]

Поэтому коэффициент корреляции обычно не определяют при нулевой
дисперсии.

\subsubsection{2.8. Ковариационная и корреляционная
матрицы}\label{ux43aux43eux432ux430ux440ux438ux430ux446ux438ux43eux43dux43dux430ux44f-ux438-ux43aux43eux440ux440ux435ux43bux44fux446ux438ux43eux43dux43dux430ux44f-ux43cux430ux442ux440ux438ux446ux44b}

Для случайного вектора

\[
X=(X_1,\ldots,X_n)^T
\]

с конечными вторыми моментами ковариационная матрица

\[
\Sigma=(\Sigma_{ij}),
\qquad
\Sigma_{ij}=\operatorname{Cov}(X_i,X_j)
\]

является матрицей Грама для центрированных векторов

\[
X_1^0,\ldots,X_n^0
\]

в \(L^2\):

\[
\Sigma_{ij}=\langle X_i^0,X_j^0\rangle.
\]

Отсюда сразу следует ее симметричность и неотрицательная определенность.
Для любого \(a=(a_1,\ldots,a_n)^T\):

\[
a^T\Sigma a
=
\sum_{i,j}a_i a_j\operatorname{Cov}(X_i,X_j)
=
\operatorname{Var}\left(\sum_{i=1}^n a_iX_i\right)
\ge0.
\]

Корреляционная матрица получается нормировкой ковариаций:

\[
\rho_{ij}
=
\rho(X_i,X_j)
=
\frac{\Sigma_{ij}}{\sqrt{\Sigma_{ii}\Sigma_{jj}}},
\]

если соответствующие дисперсии положительны. На диагонали такой матрицы
стоят единицы:

\[
R_{ii}=1.
\]

Геометрически \(R_{ij}\) - косинус угла между \(X_i^0\) и \(X_j^0\).

\subsubsection{2.9. Связь со случайными
процессами}\label{ux441ux432ux44fux437ux44c-ux441ux43e-ux441ux43bux443ux447ux430ux439ux43dux44bux43cux438-ux43fux440ux43eux446ux435ux441ux441ux430ux43cux438}

В конце ответа полезно связать билет с процессами. Для случайной
последовательности или процесса ковариации между значениями в разные
моменты времени образуют корреляционную функцию:

\[
R(s,t)=\operatorname{Cov}(X_s,X_t).
\]

В стационарном широком смысле она зависит только от лага:

\[
R(k)=\operatorname{Cov}(X_0,X_k).
\]

Это уже связь с Билет 17, Билет 24 и Билет 25: корреляционная функция и
спектральные характеристики являются развитием той же геометрии второго
порядка.

\subsection{3. Основные
определения}\label{ux43eux441ux43dux43eux432ux43dux44bux435-ux43eux43fux440ux435ux434ux435ux43bux435ux43dux438ux44f}

\begin{itemize}
\tightlist
\item
  \textbf{Центрированная случайная величина}:
\end{itemize}

\[
X^0=X-\mathbb EX.
\]

\begin{itemize}
\tightlist
\item
  \textbf{Дисперсия}:
\end{itemize}

\[
\operatorname{Var}X=\mathbb E[(X-\mathbb EX)^2]=\|X^0\|_2^2.
\]

\begin{itemize}
\tightlist
\item
  \textbf{Ковариация}:
\end{itemize}

\[
\operatorname{Cov}(X,Y)=\mathbb E[(X-\mathbb EX)(Y-\mathbb EY)].
\]

\begin{itemize}
\tightlist
\item
  \textbf{Коэффициент корреляции}:
\end{itemize}

\[
\rho(X,Y)
=
\frac{\operatorname{Cov}(X,Y)}
{\sqrt{\operatorname{Var}X\,\operatorname{Var}Y}}.
\]

\begin{itemize}
\tightlist
\item
  \textbf{Ортогональность в \(L^2\)}:
\end{itemize}

\[
X\perp Y
\quad\Longleftrightarrow\quad
\mathbb E[XY]=0.
\]

\begin{itemize}
\tightlist
\item
  \textbf{Некоррелированность}:
\end{itemize}

\[
X^0\perp Y^0
\quad\Longleftrightarrow\quad
\operatorname{Cov}(X,Y)=0.
\]

\begin{itemize}
\tightlist
\item
  \textbf{Матрица Грама}: матрица скалярных произведений векторов:
\end{itemize}

\[
G_{ij}=\langle u_i,u_j\rangle.
\]

Ковариационная матрица - матрица Грама для центрированных величин.

\subsection{4. Основные теоремы и
свойства}\label{ux43eux441ux43dux43eux432ux43dux44bux435-ux442ux435ux43eux440ux435ux43cux44b-ux438-ux441ux432ux43eux439ux441ux442ux432ux430}

\begin{itemize}
\tightlist
\item
  \(\operatorname{Cov}(X,Y)=\langle X^0,Y^0\rangle\).
\item
  \(\operatorname{Var}X=\|X^0\|_2^2\).
\item
  \(\rho(X,Y)=\cos\angle(X^0,Y^0)\) при положительных дисперсиях.
\item
  \(|\rho(X,Y)|\le1\).
\item
  \(\rho(X,Y)=0\) означает некоррелированность, то есть ортогональность
  центрированных величин.
\item
  Независимость при наличии вторых моментов влечет некоррелированность,
  но не наоборот.
\item
  \(|\rho(X,Y)|=1\) тогда и только тогда, когда \(X\) и \(Y\) линейно
  связаны почти наверное:
\end{itemize}

\[
Y=aX+b.
\]

\begin{itemize}
\tightlist
\item
  Ковариационная матрица симметрична и неотрицательно определена.
\item
  Ковариационная матрица является матрицей Грама центрированных
  случайных величин.
\end{itemize}

\subsection{5. Примеры}\label{ux43fux440ux438ux43cux435ux440ux44b}

\subsubsection{Независимые
величины}\label{ux43dux435ux437ux430ux432ux438ux441ux438ux43cux44bux435-ux432ux435ux43bux438ux447ux438ux43dux44b}

Если \(X\) и \(Y\) независимы и имеют конечные вторые моменты, то

\[
\mathbb E[XY]=\mathbb EX\,\mathbb EY.
\]

Тогда

\[
\operatorname{Cov}(X,Y)
=
\mathbb E[XY]-\mathbb EX\,\mathbb EY
=0.
\]

Значит независимые величины некоррелированы. Если дисперсии
положительны, то \(\rho(X,Y)=0\).

\subsubsection{Линейная
зависимость}\label{ux43bux438ux43dux435ux439ux43dux430ux44f-ux437ux430ux432ux438ux441ux438ux43cux43eux441ux442ux44c}

Пусть

\[
Y=aX+b,
\qquad a\ne0.
\]

Тогда

\[
Y-\mathbb EY
=
a(X-\mathbb EX).
\]

Поэтому

\[
\operatorname{Cov}(X,Y)=a\operatorname{Var}X,
\qquad
\operatorname{Var}Y=a^2\operatorname{Var}X.
\]

Если \(\operatorname{Var}X>0\), то

\[
\rho(X,Y)
=
\frac{a\operatorname{Var}X}{|a|\operatorname{Var}X}
=
\operatorname{sign}a.
\]

То есть при \(a>0\) корреляция равна \(1\), а при \(a<0\) равна \(-1\).

\subsubsection{Некоррелированные, но зависимые
величины}\label{ux43dux435ux43aux43eux440ux440ux435ux43bux438ux440ux43eux432ux430ux43dux43dux44bux435-ux43dux43e-ux437ux430ux432ux438ux441ux438ux43cux44bux435-ux432ux435ux43bux438ux447ux438ux43dux44b}

Пусть \(X\sim U[-1,1]\), а \(Y=X^2\). Тогда \(Y\) зависит от \(X\), но

\[
\operatorname{Cov}(X,Y)
=
\operatorname{Cov}(X,X^2)
=
\mathbb E X^3-\mathbb EX\,\mathbb EX^2
=0.
\]

Значит \(\rho(X,Y)=0\), если \(\operatorname{Var}Y>0\). Этот пример
показывает главную ловушку: нулевая корреляция не равна независимости.

\subsubsection{Ортогональные векторы в конечномерном
случае}\label{ux43eux440ux442ux43eux433ux43eux43dux430ux43bux44cux43dux44bux435-ux432ux435ux43aux442ux43eux440ux44b-ux432-ux43aux43eux43dux435ux447ux43dux43eux43cux435ux440ux43dux43eux43c-ux441ux43bux443ux447ux430ux435}

Если \(X=(X_1,X_2)\) - центрированный случайный вектор и

\[
\operatorname{Cov}(X_1,X_2)=0,
\]

то координаты \(X_1\) и \(X_2\) ортогональны как элементы \(L^2\). Их
ковариационная матрица диагональна:

\[
\Sigma
=
\begin{pmatrix}
\operatorname{Var}X_1 & 0\\
0 & \operatorname{Var}X_2
\end{pmatrix}.
\]

Но диагональность ковариационной матрицы сама по себе не означает
независимость координат, кроме специальных классов распределений,
например совместно гауссовского случая.

\subsection{6. Схемы доказательств /
обоснований}\label{ux441ux445ux435ux43cux44b-ux434ux43eux43aux430ux437ux430ux442ux435ux43bux44cux441ux442ux432-ux43eux431ux43eux441ux43dux43eux432ux430ux43dux438ux439}

\subsubsection{Почему ковариация - скалярное
произведение}\label{ux43fux43eux447ux435ux43cux443-ux43aux43eux432ux430ux440ux438ux430ux446ux438ux44f---ux441ux43aux430ux43bux44fux440ux43dux43eux435-ux43fux440ux43eux438ux437ux432ux435ux434ux435ux43dux438ux435}

Берем центрированные величины

\[
X^0=X-\mathbb EX,
\qquad
Y^0=Y-\mathbb EY.
\]

Тогда по определению скалярного произведения в \(L^2\):

\[
\langle X^0,Y^0\rangle
=
\mathbb E[X^0Y^0]
=
\mathbb E[(X-\mathbb EX)(Y-\mathbb EY)]
=
\operatorname{Cov}(X,Y).
\]

\subsubsection{\texorpdfstring{Почему
\(|\rho|\le1\)}{Почему \textbar\textbackslash rho\textbar\textbackslash le1}}\label{ux43fux43eux447ux435ux43cux443-rhole1}

Применяем Коши-Буняковского в \(L^2\):

\[
|\langle X^0,Y^0\rangle|
\le
\|X^0\|_2\|Y^0\|_2.
\]

После подстановки получаем

\[
|\operatorname{Cov}(X,Y)|
\le
\sqrt{\operatorname{Var}X\,\operatorname{Var}Y}.
\]

Если дисперсии положительны, делим на знаменатель и получаем

\[
|\rho(X,Y)|\le1.
\]

\subsubsection{\texorpdfstring{Почему \(|\rho|=1\) означает линейную
зависимость}{Почему \textbar\textbackslash rho\textbar=1 означает линейную зависимость}}\label{ux43fux43eux447ux435ux43cux443-rho1-ux43eux437ux43dux430ux447ux430ux435ux442-ux43bux438ux43dux435ux439ux43dux443ux44e-ux437ux430ux432ux438ux441ux438ux43cux43eux441ux442ux44c}

В неравенстве Коши-Буняковского равенство достигается тогда и только
тогда, когда векторы линейно зависимы. Поэтому

\[
|\rho(X,Y)|=1
\]

эквивалентно существованию числа \(a\ne0\) такого, что

\[
Y^0=aX^0
\quad\text{почти наверное}.
\]

Раскрывая центрирование:

\[
Y-\mathbb EY=a(X-\mathbb EX),
\]

то есть

\[
Y=aX+(\mathbb EY-a\mathbb EX)
\quad\text{почти наверное}.
\]

\subsubsection{Почему ковариационная матрица неотрицательно
определена}\label{ux43fux43eux447ux435ux43cux443-ux43aux43eux432ux430ux440ux438ux430ux446ux438ux43eux43dux43dux430ux44f-ux43cux430ux442ux440ux438ux446ux430-ux43dux435ux43eux442ux440ux438ux446ux430ux442ux435ux43bux44cux43dux43e-ux43eux43fux440ux435ux434ux435ux43bux435ux43dux430}

Пусть

\[
\Sigma_{ij}=\operatorname{Cov}(X_i,X_j).
\]

Тогда для любого \(a\in\mathbb R^n\):

\[
a^T\Sigma a
=
\sum_{i,j}a_i a_j\operatorname{Cov}(X_i,X_j)
=
\operatorname{Var}\left(\sum_i a_iX_i\right)
\ge0.
\]

Значит \(\Sigma\) неотрицательно определена.

\subsection{7. Что написать на
доске}\label{ux447ux442ux43e-ux43dux430ux43fux438ux441ux430ux442ux44c-ux43dux430-ux434ux43eux441ux43aux435}

\begin{enumerate}
\def\labelenumi{\arabic{enumi}.}
\tightlist
\item
  Центрирование:
\end{enumerate}

\[
X^0=X-\mathbb EX.
\]

\begin{enumerate}
\def\labelenumi{\arabic{enumi}.}
\setcounter{enumi}{1}
\tightlist
\item
  Ковариация как скалярное произведение:
\end{enumerate}

\[
\operatorname{Cov}(X,Y)=\langle X^0,Y^0\rangle.
\]

\begin{enumerate}
\def\labelenumi{\arabic{enumi}.}
\setcounter{enumi}{2}
\tightlist
\item
  Дисперсия как квадрат длины:
\end{enumerate}

\[
\operatorname{Var}X=\|X^0\|_2^2.
\]

\begin{enumerate}
\def\labelenumi{\arabic{enumi}.}
\setcounter{enumi}{3}
\tightlist
\item
  Коэффициент корреляции:
\end{enumerate}

\[
\rho(X,Y)
=
\frac{\operatorname{Cov}(X,Y)}
{\sqrt{\operatorname{Var}X\,\operatorname{Var}Y}}
=
\frac{\langle X^0,Y^0\rangle}{\|X^0\|_2\|Y^0\|_2}
=
\cos\angle(X^0,Y^0).
\]

\begin{enumerate}
\def\labelenumi{\arabic{enumi}.}
\setcounter{enumi}{4}
\tightlist
\item
  Главное следствие:
\end{enumerate}

\[
|\rho(X,Y)|\le1.
\]

\begin{enumerate}
\def\labelenumi{\arabic{enumi}.}
\setcounter{enumi}{5}
\tightlist
\item
  Крайние случаи:
\end{enumerate}

\[
\rho=0 \Longleftrightarrow X^0\perp Y^0,
\]

\[
|\rho|=1 \Longleftrightarrow Y=aX+b\quad\text{п.н.}
\]

\subsection{8. Типичные дополнительные
вопросы}\label{ux442ux438ux43fux438ux447ux43dux44bux435-ux434ux43eux43fux43eux43bux43dux438ux442ux435ux43bux44cux43dux44bux435-ux432ux43eux43fux440ux43eux441ux44b}

\textbf{Почему нужно центрировать величины?}

Потому что корреляция должна измерять совместные отклонения от средних,
а не сами уровни величин. Без центрирования скалярное произведение
\(\mathbb E[XY]\) зависит от сдвигов. Например, если заменить \(X\) на
\(X+c\), то \(\mathbb E[(X+c)Y]\) обычно изменится, хотя форма
случайного отклонения \(X-\mathbb EX\) не изменилась. Центрирование
убирает постоянную часть:

\[
X^0=X-\mathbb EX.
\]

\textbf{Почему коэффициент корреляции не определяют при нулевой
дисперсии?}

Если \(\operatorname{Var}X=0\), то \(X=\mathbb EX\) почти наверное и

\[
X^0=0.
\]

У нулевого вектора нет направления и угла с другим вектором. В формуле
для \(\rho\) появляется деление на ноль:

\[
\rho(X,Y)=
\frac{\operatorname{Cov}(X,Y)}
{\sqrt{\operatorname{Var}X\,\operatorname{Var}Y}}.
\]

\textbf{В чем разница между ковариацией и корреляцией?}

Ковариация - это ненормированное скалярное произведение центрированных
величин:

\[
\operatorname{Cov}(X,Y)=\langle X^0,Y^0\rangle.
\]

Она зависит от масштаба измерения. Если \(X\) умножить на 10, ковариация
тоже умножится на 10. Корреляция - нормированная ковариация:

\[
\rho(X,Y)=
\frac{\operatorname{Cov}(X,Y)}
{\sqrt{\operatorname{Var}X\,\operatorname{Var}Y}},
\]

поэтому всегда лежит в \([-1,1]\) и показывает только направление
линейной связи.

\textbf{Почему \(|\rho|\le1\)?}

Потому что в \(L^2\) выполняется неравенство Коши-Буняковского:

\[
|\langle X^0,Y^0\rangle|
\le
\|X^0\|_2\|Y^0\|_2.
\]

После подстановки \(\langle X^0,Y^0\rangle=\operatorname{Cov}(X,Y)\) и
\(\|X^0\|_2^2=\operatorname{Var}X\) получаем

\[
|\operatorname{Cov}(X,Y)|
\le
\sqrt{\operatorname{Var}X\,\operatorname{Var}Y},
\]

а значит \(|\rho|\le1\).

\textbf{Что означает \(\rho=0\) геометрически?}

Это означает, что центрированные величины ортогональны в \(L^2\):

\[
X^0\perp Y^0.
\]

Вероятностно это называется некоррелированностью:

\[
\operatorname{Cov}(X,Y)=0.
\]

\textbf{Следует ли независимость из \(\rho=0\)?}

В общем случае нет. Нулевая корреляция означает только отсутствие
линейной связи второго порядка. Например, если \(X\sim U[-1,1]\) и
\(Y=X^2\), то \(Y\) зависит от \(X\), но

\[
\operatorname{Cov}(X,Y)=\operatorname{Cov}(X,X^2)=0.
\]

\textbf{Когда из некоррелированности все-таки следует независимость?}

Важный случай - совместно гауссовские величины. Если вектор \((X,Y)\)
гауссовский и

\[
\operatorname{Cov}(X,Y)=0,
\]

то \(X\) и \(Y\) независимы. Это специальное свойство гауссовских
распределений, а не общее правило.

\textbf{Что означает \(|\rho|=1\)?}

Это означает, что в неравенстве Коши-Буняковского достигнуто равенство,
поэтому центрированные величины линейно зависимы:

\[
Y^0=aX^0
\quad\text{п.н.}
\]

Эквивалентно,

\[
Y=aX+b
\quad\text{п.н.}
\]

При \(a>0\) получаем \(\rho=1\), при \(a<0\) получаем \(\rho=-1\).

\textbf{Почему ковариационная матрица неотрицательно определена?}

Пусть \(\Sigma_{ij}=\operatorname{Cov}(X_i,X_j)\). Тогда для любого
\(a\in\mathbb R^n\):

\[
a^T\Sigma a
=
\operatorname{Var}\left(\sum_i a_iX_i\right)
\ge0.
\]

Значит \(\Sigma\) неотрицательно определена. Геометрически это матрица
Грама центрированных векторов \(X_1^0,\ldots,X_n^0\).

\textbf{Как билет связан с Билет 12, Билет 13, Билет 18, Билет 24 и
Билет 25?}

С Билетом 12 связь через гильбертово пространство \(L^2\). С Билетом 13
- через получение моментов и ковариационной матрицы из
характеристической функции. С Билетом 18 - через гауссовские векторы,
где некоррелированность равносильна независимости. С Билетом 24 и
Билетом 25 - через корреляционные функции и спектральные характеристики
случайных последовательностей.

\subsection{9. Типичные
ошибки}\label{ux442ux438ux43fux438ux447ux43dux44bux435-ux43eux448ux438ux431ux43aux438}

\begin{itemize}
\tightlist
\item
  Делать вывод о независимости из \(\rho=0\).
\item
  Забывать центрирование и писать \(\mathbb E[XY]\) вместо
  \(\operatorname{Cov}(X,Y)\) без условия \(\mathbb EX=\mathbb EY=0\).
\item
  Не проверять существование вторых моментов.
\item
  Не проверять, что дисперсии положительны.
\item
  Называть ковариационную матрицу корреляционной без пояснения
  нормировки.
\item
  Путать нулевую ковариацию и ортогональность самих \(X,Y\): на самом
  деле ортогональны центрированные величины \(X^0,Y^0\).
\item
  Забывать, что элементы \(L^2\) рассматриваются с точностью до
  равенства почти наверное.
\end{itemize}

\subsection{10. Связи с другими
билетами}\label{ux441ux432ux44fux437ux438-ux441-ux434ux440ux443ux433ux438ux43cux438-ux431ux438ux43bux435ux442ux430ux43cux438}

\begin{itemize}
\tightlist
\item
  Билет 11: математические ожидания, дисперсии и ковариации выражаются
  через интегралы.
\item
  Билет 12: \(L^2\) - гильбертово пространство; корреляция использует
  его скалярное произведение.
\item
  Билет 13: моменты и ковариационная матрица могут быть получены из
  производных характеристической функции.
\item
  Билет 18: для совместно гауссовских величин некоррелированность
  эквивалентна независимости.
\item
  Билет 24: корреляционная функция процесса - ковариация значений
  процесса в разные моменты времени.
\item
  Билет 25: спектральные характеристики строятся из корреляционной
  функции стационарной последовательности.
\end{itemize}

\subsection{11. Минимум на
удовлетворительно}\label{ux43cux438ux43dux438ux43cux443ux43c-ux43dux430-ux443ux434ux43eux432ux43bux435ux442ux432ux43eux440ux438ux442ux435ux43bux44cux43dux43e}

Нужно сказать:

\begin{enumerate}
\def\labelenumi{\arabic{enumi}.}
\tightlist
\item
  Для \(X,Y\in L^2\) определены центрированные величины:
\end{enumerate}

\[
X^0=X-\mathbb EX,
\qquad
Y^0=Y-\mathbb EY.
\]

\begin{enumerate}
\def\labelenumi{\arabic{enumi}.}
\setcounter{enumi}{1}
\tightlist
\item
  Ковариация - скалярное произведение центрированных величин:
\end{enumerate}

\[
\operatorname{Cov}(X,Y)=\langle X^0,Y^0\rangle.
\]

\begin{enumerate}
\def\labelenumi{\arabic{enumi}.}
\setcounter{enumi}{2}
\tightlist
\item
  Коэффициент корреляции:
\end{enumerate}

\[
\rho(X,Y)=
\frac{\operatorname{Cov}(X,Y)}
{\sqrt{\operatorname{Var}X\,\operatorname{Var}Y}}.
\]

\begin{enumerate}
\def\labelenumi{\arabic{enumi}.}
\setcounter{enumi}{3}
\tightlist
\item
  Геометрический смысл:
\end{enumerate}

\[
\rho(X,Y)=\cos\angle(X^0,Y^0),
\qquad
|\rho|\le1.
\]

\begin{enumerate}
\def\labelenumi{\arabic{enumi}.}
\setcounter{enumi}{4}
\tightlist
\item
  Главная ловушка: \(\rho=0\) не означает независимость.
\end{enumerate}

\subsection{12. Минимум на
хорошо}\label{ux43cux438ux43dux438ux43cux443ux43c-ux43dux430-ux445ux43eux440ux43eux448ux43e}

Помимо минимума, нужно объяснить:

\begin{itemize}
\tightlist
\item
  почему нужно условие \(\mathbb EX^2,\mathbb EY^2<\infty\);
\item
  почему нужны положительные дисперсии для \(\rho\);
\item
  как из Коши-Буняковского получается \(|\rho|\le1\);
\item
  что \(\rho=0\) означает ортогональность центрированных величин;
\item
  что \(|\rho|=1\) означает линейную зависимость почти наверное;
\item
  привести пример зависимых, но некоррелированных величин:
  \(X\sim U[-1,1]\), \(Y=X^2\).
\end{itemize}

\subsection{13. Что добавить для
отлично}\label{ux447ux442ux43e-ux434ux43eux431ux430ux432ux438ux442ux44c-ux434ux43bux44f-ux43eux442ux43bux438ux447ux43dux43e}

Для сильного ответа стоит добавить:

\begin{itemize}
\tightlist
\item
  формулировку через гильбертово пространство
  \(L^2(\Omega,\mathcal F,P)\);
\item
  замечание, что элементы \(L^2\) - классы эквивалентности по равенству
  почти наверное;
\item
  доказательство \(|\rho|\le1\) через Коши-Буняковского;
\item
  критерий равенства \(|\rho|=1\);
\item
  объяснение ковариационной матрицы как матрицы Грама;
\item
  доказательство неотрицательной определенности ковариационной матрицы;
\item
  связь с гауссовским случаем: для совместно гауссовских величин
  некоррелированность дает независимость;
\item
  связь с корреляционной функцией случайных процессов и спектральными
  характеристиками.
\end{itemize}

\subsection{14. Устная версия ответа на 5
минут}\label{ux443ux441ux442ux43dux430ux44f-ux432ux435ux440ux441ux438ux44f-ux43eux442ux432ux435ux442ux430-ux43dux430-5-ux43cux438ux43dux443ux442}

\begin{enumerate}
\def\labelenumi{\arabic{enumi}.}
\tightlist
\item
  В этом билете случайные величины с конечным вторым моментом
  рассматриваются как векторы в \(L^2\).
\item
  В \(L^2\) скалярное произведение:
\end{enumerate}

\[
\langle X,Y\rangle=\mathbb E[XY].
\]

\begin{enumerate}
\def\labelenumi{\arabic{enumi}.}
\setcounter{enumi}{2}
\tightlist
\item
  Для геометрии корреляции надо перейти к центрированным величинам:
\end{enumerate}

\[
X^0=X-\mathbb EX,
\qquad
Y^0=Y-\mathbb EY.
\]

\begin{enumerate}
\def\labelenumi{\arabic{enumi}.}
\setcounter{enumi}{3}
\tightlist
\item
  Тогда
\end{enumerate}

\[
\operatorname{Cov}(X,Y)=\langle X^0,Y^0\rangle,
\qquad
\operatorname{Var}X=\|X^0\|_2^2.
\]

\begin{enumerate}
\def\labelenumi{\arabic{enumi}.}
\setcounter{enumi}{4}
\tightlist
\item
  При положительных дисперсиях коэффициент корреляции:
\end{enumerate}

\[
\rho(X,Y)=
\frac{\operatorname{Cov}(X,Y)}
{\sqrt{\operatorname{Var}X\,\operatorname{Var}Y}}
=
\frac{\langle X^0,Y^0\rangle}{\|X^0\|_2\|Y^0\|_2}.
\]

\begin{enumerate}
\def\labelenumi{\arabic{enumi}.}
\setcounter{enumi}{5}
\tightlist
\item
  Значит \(\rho\) - косинус угла между \(X^0\) и \(Y^0\):
\end{enumerate}

\[
\rho(X,Y)=\cos\angle(X^0,Y^0).
\]

\begin{enumerate}
\def\labelenumi{\arabic{enumi}.}
\setcounter{enumi}{6}
\tightlist
\item
  По Коши-Буняковскому:
\end{enumerate}

\[
|\rho(X,Y)|\le1.
\]

\begin{enumerate}
\def\labelenumi{\arabic{enumi}.}
\setcounter{enumi}{7}
\tightlist
\item
  Если \(\rho=0\), то \(X^0\) и \(Y^0\) ортогональны, то есть величины
  некоррелированы. Это не значит, что они независимы.
\item
  Если \(|\rho|=1\), то центрированные величины линейно зависимы, то
  есть \(Y=aX+b\) почти наверное.
\item
  Пример ловушки: \(X\sim U[-1,1]\), \(Y=X^2\). Они зависимы, но
  \(\operatorname{Cov}(X,Y)=0\).
\item
  Для случайного вектора ковариационная матрица - матрица Грама
  центрированных координат, поэтому она симметрична и неотрицательно
  определена.
\item
  Завершить можно связью: эта геометрия лежит в основе корреляционных
  функций процессов и спектральных характеристик.
\end{enumerate}

\subsection{15. Если осталось 2 минуты перед
экзаменом}\label{ux435ux441ux43bux438-ux43eux441ux442ux430ux43bux43eux441ux44c-2-ux43cux438ux43dux443ux442ux44b-ux43fux435ux440ux435ux434-ux44dux43aux437ux430ux43cux435ux43dux43eux43c}

Запомнить три формулы:

\[
X^0=X-\mathbb EX,
\]

\[
\operatorname{Cov}(X,Y)=\langle X^0,Y^0\rangle,
\]

\[
\rho(X,Y)=
\frac{\operatorname{Cov}(X,Y)}
{\sqrt{\operatorname{Var}X\,\operatorname{Var}Y}}
=
\cos\angle(X^0,Y^0).
\]

Главное следствие:

\[
|\rho|\le1.
\]

Главная ловушка: \(\rho=0\) означает только некоррелированность, а не
независимость. Независимость из некоррелированности следует не всегда;
важное исключение - совместно гауссовские величины.

\newpage

\end{document}
